\captionsetup{format=hang, justification=raggedright} %Formatierung der Bildunterschriften
\setcounter{secnumdepth}{4} %Tiefe der Gliederung
\setcounter{tocdepth}{4}   %Tiefe der Gliederung im Inhaltsverzeichnis

\counterwithout{table}{section}

\addbibresource{References/bibliography.bib} %BibLatex-Datenbank
\setlength{\bibitemsep}{5pt plus 0.3ex} %zum setzen von Abständen zwischen den Quellen im Literaturverzeichnis

%TODO Zitat Platzhalter, Verwendung: \citeTODO{https://google.com}
\newcommand{\citeTODO}[1]{%
    \textbf{[ZITAT]}%
}

%Source für mehrere Quellen in einem Zitat wie. \citeTODO{\source{https://google.com}, {https://google.com}}
\newcommand{\source}[1]{#1}

%TODO Platzhalter, Verwendung: \TODO{noch zu machen}
\newcommand{\TODO}[1]{%
    \textbf{[TODO]}%
}

\DeclareCiteCommand{\pagecite}
  {}
  {\iffieldundef{postnote}{\PackageError{pagecite}{No page number given}}{\printtext[bibhyperref]{\mkbibparens{\printfield{postnote}}}}
  }
  {}
  {}

% This command allows page citations like \pagecite[20] to produce (p.20)
% and use the last cited key with \textcite
% Attention: If no correct prior \textcite is found, an error is raised or
% it will use the last one.
% ------------------------------------------------------------------

%% Store the last cite key used with \textcite
%\newcommand{\lastcitekey}{}
%
%% Wrapper for \textcite to store the last cite key used
%\let\oldtextcite\textcite
%\renewcommand{\textcite}[1]{%
%    \renewcommand{\lastcitekey}{#1}%
%    \oldtextcite{#1}%
%}
%
%\DeclareCiteCommand{\pagecitecommand}[\mkbibparens]
%  {}
%  {\printtext[bibhyperref]{\printfield{postnote}}}
%  {}
%  {}
%
%\newcommand{\pagecite}[1]{%
%    \ifx\lastcitekey\empty
%        \PackageError{pagecite}
%    \else
%        \pagecitecommand[#1]{\lastcitekey}%
%    \fi
%}
% ------------------------------------------------------------------